\documentclass{article}

\usepackage{iclr2026_conference,times}

%%%%% NEW MATH DEFINITIONS %%%%%

\usepackage{amsmath,amsfonts,bm}

% Mark sections of captions for referring to divisions of figures
\newcommand{\figleft}{{\em (Left)}}
\newcommand{\figcenter}{{\em (Center)}}
\newcommand{\figright}{{\em (Right)}}
\newcommand{\figtop}{{\em (Top)}}
\newcommand{\figbottom}{{\em (Bottom)}}
\newcommand{\captiona}{{\em (a)}}
\newcommand{\captionb}{{\em (b)}}
\newcommand{\captionc}{{\em (c)}}
\newcommand{\captiond}{{\em (d)}}

% Highlight a newly defined term
\newcommand{\newterm}[1]{{\bf #1}}


% Figure reference, lower-case.
\def\figref#1{figure~\ref{#1}}
% Figure reference, capital. For start of sentence
\def\Figref#1{Figure~\ref{#1}}
\def\twofigref#1#2{figures \ref{#1} and \ref{#2}}
\def\quadfigref#1#2#3#4{figures \ref{#1}, \ref{#2}, \ref{#3} and \ref{#4}}
% Section reference, lower-case.
\def\secref#1{section~\ref{#1}}
% Section reference, capital.
\def\Secref#1{Section~\ref{#1}}
% Reference to two sections.
\def\twosecrefs#1#2{sections \ref{#1} and \ref{#2}}
% Reference to three sections.
\def\secrefs#1#2#3{sections \ref{#1}, \ref{#2} and \ref{#3}}
% Reference to an equation, lower-case.
\def\eqref#1{equation~\ref{#1}}
% Reference to an equation, upper case
\def\Eqref#1{Equation~\ref{#1}}
% A raw reference to an equation---avoid using if possible
\def\plaineqref#1{\ref{#1}}
% Reference to a chapter, lower-case.
\def\chapref#1{chapter~\ref{#1}}
% Reference to an equation, upper case.
\def\Chapref#1{Chapter~\ref{#1}}
% Reference to a range of chapters
\def\rangechapref#1#2{chapters\ref{#1}--\ref{#2}}
% Reference to an algorithm, lower-case.
\def\algref#1{algorithm~\ref{#1}}
% Reference to an algorithm, upper case.
\def\Algref#1{Algorithm~\ref{#1}}
\def\twoalgref#1#2{algorithms \ref{#1} and \ref{#2}}
\def\Twoalgref#1#2{Algorithms \ref{#1} and \ref{#2}}
% Reference to a part, lower case
\def\partref#1{part~\ref{#1}}
% Reference to a part, upper case
\def\Partref#1{Part~\ref{#1}}
\def\twopartref#1#2{parts \ref{#1} and \ref{#2}}

\def\ceil#1{\lceil #1 \rceil}
\def\floor#1{\lfloor #1 \rfloor}
\def\1{\bm{1}}
\newcommand{\train}{\mathcal{D}}
\newcommand{\valid}{\mathcal{D_{\mathrm{valid}}}}
\newcommand{\test}{\mathcal{D_{\mathrm{test}}}}

\def\eps{{\epsilon}}


% Random variables
\def\reta{{\textnormal{$\eta$}}}
\def\ra{{\textnormal{a}}}
\def\rb{{\textnormal{b}}}
\def\rc{{\textnormal{c}}}
\def\rd{{\textnormal{d}}}
\def\re{{\textnormal{e}}}
\def\rf{{\textnormal{f}}}
\def\rg{{\textnormal{g}}}
\def\rh{{\textnormal{h}}}
\def\ri{{\textnormal{i}}}
\def\rj{{\textnormal{j}}}
\def\rk{{\textnormal{k}}}
\def\rl{{\textnormal{l}}}
% rm is already a command, just don't name any random variables m
\def\rn{{\textnormal{n}}}
\def\ro{{\textnormal{o}}}
\def\rp{{\textnormal{p}}}
\def\rq{{\textnormal{q}}}
\def\rr{{\textnormal{r}}}
\def\rs{{\textnormal{s}}}
\def\rt{{\textnormal{t}}}
\def\ru{{\textnormal{u}}}
\def\rv{{\textnormal{v}}}
\def\rw{{\textnormal{w}}}
\def\rx{{\textnormal{x}}}
\def\ry{{\textnormal{y}}}
\def\rz{{\textnormal{z}}}

% Random vectors
\def\rvepsilon{{\mathbf{\epsilon}}}
\def\rvtheta{{\mathbf{\theta}}}
\def\rva{{\mathbf{a}}}
\def\rvb{{\mathbf{b}}}
\def\rvc{{\mathbf{c}}}
\def\rvd{{\mathbf{d}}}
\def\rve{{\mathbf{e}}}
\def\rvf{{\mathbf{f}}}
\def\rvg{{\mathbf{g}}}
\def\rvh{{\mathbf{h}}}
\def\rvu{{\mathbf{i}}}
\def\rvj{{\mathbf{j}}}
\def\rvk{{\mathbf{k}}}
\def\rvl{{\mathbf{l}}}
\def\rvm{{\mathbf{m}}}
\def\rvn{{\mathbf{n}}}
\def\rvo{{\mathbf{o}}}
\def\rvp{{\mathbf{p}}}
\def\rvq{{\mathbf{q}}}
\def\rvr{{\mathbf{r}}}
\def\rvs{{\mathbf{s}}}
\def\rvt{{\mathbf{t}}}
\def\rvu{{\mathbf{u}}}
\def\rvv{{\mathbf{v}}}
\def\rvw{{\mathbf{w}}}
\def\rvx{{\mathbf{x}}}
\def\rvy{{\mathbf{y}}}
\def\rvz{{\mathbf{z}}}

% Elements of random vectors
\def\erva{{\textnormal{a}}}
\def\ervb{{\textnormal{b}}}
\def\ervc{{\textnormal{c}}}
\def\ervd{{\textnormal{d}}}
\def\erve{{\textnormal{e}}}
\def\ervf{{\textnormal{f}}}
\def\ervg{{\textnormal{g}}}
\def\ervh{{\textnormal{h}}}
\def\ervi{{\textnormal{i}}}
\def\ervj{{\textnormal{j}}}
\def\ervk{{\textnormal{k}}}
\def\ervl{{\textnormal{l}}}
\def\ervm{{\textnormal{m}}}
\def\ervn{{\textnormal{n}}}
\def\ervo{{\textnormal{o}}}
\def\ervp{{\textnormal{p}}}
\def\ervq{{\textnormal{q}}}
\def\ervr{{\textnormal{r}}}
\def\ervs{{\textnormal{s}}}
\def\ervt{{\textnormal{t}}}
\def\ervu{{\textnormal{u}}}
\def\ervv{{\textnormal{v}}}
\def\ervw{{\textnormal{w}}}
\def\ervx{{\textnormal{x}}}
\def\ervy{{\textnormal{y}}}
\def\ervz{{\textnormal{z}}}

% Random matrices
\def\rmA{{\mathbf{A}}}
\def\rmB{{\mathbf{B}}}
\def\rmC{{\mathbf{C}}}
\def\rmD{{\mathbf{D}}}
\def\rmE{{\mathbf{E}}}
\def\rmF{{\mathbf{F}}}
\def\rmG{{\mathbf{G}}}
\def\rmH{{\mathbf{H}}}
\def\rmI{{\mathbf{I}}}
\def\rmJ{{\mathbf{J}}}
\def\rmK{{\mathbf{K}}}
\def\rmL{{\mathbf{L}}}
\def\rmM{{\mathbf{M}}}
\def\rmN{{\mathbf{N}}}
\def\rmO{{\mathbf{O}}}
\def\rmP{{\mathbf{P}}}
\def\rmQ{{\mathbf{Q}}}
\def\rmR{{\mathbf{R}}}
\def\rmS{{\mathbf{S}}}
\def\rmT{{\mathbf{T}}}
\def\rmU{{\mathbf{U}}}
\def\rmV{{\mathbf{V}}}
\def\rmW{{\mathbf{W}}}
\def\rmX{{\mathbf{X}}}
\def\rmY{{\mathbf{Y}}}
\def\rmZ{{\mathbf{Z}}}

% Elements of random matrices
\def\ermA{{\textnormal{A}}}
\def\ermB{{\textnormal{B}}}
\def\ermC{{\textnormal{C}}}
\def\ermD{{\textnormal{D}}}
\def\ermE{{\textnormal{E}}}
\def\ermF{{\textnormal{F}}}
\def\ermG{{\textnormal{G}}}
\def\ermH{{\textnormal{H}}}
\def\ermI{{\textnormal{I}}}
\def\ermJ{{\textnormal{J}}}
\def\ermK{{\textnormal{K}}}
\def\ermL{{\textnormal{L}}}
\def\ermM{{\textnormal{M}}}
\def\ermN{{\textnormal{N}}}
\def\ermO{{\textnormal{O}}}
\def\ermP{{\textnormal{P}}}
\def\ermQ{{\textnormal{Q}}}
\def\ermR{{\textnormal{R}}}
\def\ermS{{\textnormal{S}}}
\def\ermT{{\textnormal{T}}}
\def\ermU{{\textnormal{U}}}
\def\ermV{{\textnormal{V}}}
\def\ermW{{\textnormal{W}}}
\def\ermX{{\textnormal{X}}}
\def\ermY{{\textnormal{Y}}}
\def\ermZ{{\textnormal{Z}}}

% Vectors
\def\vzero{{\bm{0}}}
\def\vone{{\bm{1}}}
\def\vmu{{\bm{\mu}}}
\def\vtheta{{\bm{\theta}}}
\def\va{{\bm{a}}}
\def\vb{{\bm{b}}}
\def\vc{{\bm{c}}}
\def\vd{{\bm{d}}}
\def\ve{{\bm{e}}}
\def\vf{{\bm{f}}}
\def\vg{{\bm{g}}}
\def\vh{{\bm{h}}}
\def\vi{{\bm{i}}}
\def\vj{{\bm{j}}}
\def\vk{{\bm{k}}}
\def\vl{{\bm{l}}}
\def\vm{{\bm{m}}}
\def\vn{{\bm{n}}}
\def\vo{{\bm{o}}}
\def\vp{{\bm{p}}}
\def\vq{{\bm{q}}}
\def\vr{{\bm{r}}}
\def\vs{{\bm{s}}}
\def\vt{{\bm{t}}}
\def\vu{{\bm{u}}}
\def\vv{{\bm{v}}}
\def\vw{{\bm{w}}}
\def\vx{{\bm{x}}}
\def\vy{{\bm{y}}}
\def\vz{{\bm{z}}}

% Elements of vectors
\def\evalpha{{\alpha}}
\def\evbeta{{\beta}}
\def\evepsilon{{\epsilon}}
\def\evlambda{{\lambda}}
\def\evomega{{\omega}}
\def\evmu{{\mu}}
\def\evpsi{{\psi}}
\def\evsigma{{\sigma}}
\def\evtheta{{\theta}}
\def\eva{{a}}
\def\evb{{b}}
\def\evc{{c}}
\def\evd{{d}}
\def\eve{{e}}
\def\evf{{f}}
\def\evg{{g}}
\def\evh{{h}}
\def\evi{{i}}
\def\evj{{j}}
\def\evk{{k}}
\def\evl{{l}}
\def\evm{{m}}
\def\evn{{n}}
\def\evo{{o}}
\def\evp{{p}}
\def\evq{{q}}
\def\evr{{r}}
\def\evs{{s}}
\def\evt{{t}}
\def\evu{{u}}
\def\evv{{v}}
\def\evw{{w}}
\def\evx{{x}}
\def\evy{{y}}
\def\evz{{z}}

% Matrix
\def\mA{{\bm{A}}}
\def\mB{{\bm{B}}}
\def\mC{{\bm{C}}}
\def\mD{{\bm{D}}}
\def\mE{{\bm{E}}}
\def\mF{{\bm{F}}}
\def\mG{{\bm{G}}}
\def\mH{{\bm{H}}}
\def\mI{{\bm{I}}}
\def\mJ{{\bm{J}}}
\def\mK{{\bm{K}}}
\def\mL{{\bm{L}}}
\def\mM{{\bm{M}}}
\def\mN{{\bm{N}}}
\def\mO{{\bm{O}}}
\def\mP{{\bm{P}}}
\def\mQ{{\bm{Q}}}
\def\mR{{\bm{R}}}
\def\mS{{\bm{S}}}
\def\mT{{\bm{T}}}
\def\mU{{\bm{U}}}
\def\mV{{\bm{V}}}
\def\mW{{\bm{W}}}
\def\mX{{\bm{X}}}
\def\mY{{\bm{Y}}}
\def\mZ{{\bm{Z}}}
\def\mBeta{{\bm{\beta}}}
\def\mPhi{{\bm{\Phi}}}
\def\mLambda{{\bm{\Lambda}}}
\def\mSigma{{\bm{\Sigma}}}

% Tensor
\DeclareMathAlphabet{\mathsfit}{\encodingdefault}{\sfdefault}{m}{sl}
\SetMathAlphabet{\mathsfit}{bold}{\encodingdefault}{\sfdefault}{bx}{n}
\newcommand{\tens}[1]{\bm{\mathsfit{#1}}}
\def\tA{{\tens{A}}}
\def\tB{{\tens{B}}}
\def\tC{{\tens{C}}}
\def\tD{{\tens{D}}}
\def\tE{{\tens{E}}}
\def\tF{{\tens{F}}}
\def\tG{{\tens{G}}}
\def\tH{{\tens{H}}}
\def\tI{{\tens{I}}}
\def\tJ{{\tens{J}}}
\def\tK{{\tens{K}}}
\def\tL{{\tens{L}}}
\def\tM{{\tens{M}}}
\def\tN{{\tens{N}}}
\def\tO{{\tens{O}}}
\def\tP{{\tens{P}}}
\def\tQ{{\tens{Q}}}
\def\tR{{\tens{R}}}
\def\tS{{\tens{S}}}
\def\tT{{\tens{T}}}
\def\tU{{\tens{U}}}
\def\tV{{\tens{V}}}
\def\tW{{\tens{W}}}
\def\tX{{\tens{X}}}
\def\tY{{\tens{Y}}}
\def\tZ{{\tens{Z}}}


% Graph
\def\gA{{\mathcal{A}}}
\def\gB{{\mathcal{B}}}
\def\gC{{\mathcal{C}}}
\def\gD{{\mathcal{D}}}
\def\gE{{\mathcal{E}}}
\def\gF{{\mathcal{F}}}
\def\gG{{\mathcal{G}}}
\def\gH{{\mathcal{H}}}
\def\gI{{\mathcal{I}}}
\def\gJ{{\mathcal{J}}}
\def\gK{{\mathcal{K}}}
\def\gL{{\mathcal{L}}}
\def\gM{{\mathcal{M}}}
\def\gN{{\mathcal{N}}}
\def\gO{{\mathcal{O}}}
\def\gP{{\mathcal{P}}}
\def\gQ{{\mathcal{Q}}}
\def\gR{{\mathcal{R}}}
\def\gS{{\mathcal{S}}}
\def\gT{{\mathcal{T}}}
\def\gU{{\mathcal{U}}}
\def\gV{{\mathcal{V}}}
\def\gW{{\mathcal{W}}}
\def\gX{{\mathcal{X}}}
\def\gY{{\mathcal{Y}}}
\def\gZ{{\mathcal{Z}}}

% Sets
\def\sA{{\mathbb{A}}}
\def\sB{{\mathbb{B}}}
\def\sC{{\mathbb{C}}}
\def\sD{{\mathbb{D}}}
% Don't use a set called E, because this would be the same as our symbol
% for expectation.
\def\sF{{\mathbb{F}}}
\def\sG{{\mathbb{G}}}
\def\sH{{\mathbb{H}}}
\def\sI{{\mathbb{I}}}
\def\sJ{{\mathbb{J}}}
\def\sK{{\mathbb{K}}}
\def\sL{{\mathbb{L}}}
\def\sM{{\mathbb{M}}}
\def\sN{{\mathbb{N}}}
\def\sO{{\mathbb{O}}}
\def\sP{{\mathbb{P}}}
\def\sQ{{\mathbb{Q}}}
\def\sR{{\mathbb{R}}}
\def\sS{{\mathbb{S}}}
\def\sT{{\mathbb{T}}}
\def\sU{{\mathbb{U}}}
\def\sV{{\mathbb{V}}}
\def\sW{{\mathbb{W}}}
\def\sX{{\mathbb{X}}}
\def\sY{{\mathbb{Y}}}
\def\sZ{{\mathbb{Z}}}

% Entries of a matrix
\def\emLambda{{\Lambda}}
\def\emA{{A}}
\def\emB{{B}}
\def\emC{{C}}
\def\emD{{D}}
\def\emE{{E}}
\def\emF{{F}}
\def\emG{{G}}
\def\emH{{H}}
\def\emI{{I}}
\def\emJ{{J}}
\def\emK{{K}}
\def\emL{{L}}
\def\emM{{M}}
\def\emN{{N}}
\def\emO{{O}}
\def\emP{{P}}
\def\emQ{{Q}}
\def\emR{{R}}
\def\emS{{S}}
\def\emT{{T}}
\def\emU{{U}}
\def\emV{{V}}
\def\emW{{W}}
\def\emX{{X}}
\def\emY{{Y}}
\def\emZ{{Z}}
\def\emSigma{{\Sigma}}

% entries of a tensor
% Same font as tensor, without \bm wrapper
\newcommand{\etens}[1]{\mathsfit{#1}}
\def\etLambda{{\etens{\Lambda}}}
\def\etA{{\etens{A}}}
\def\etB{{\etens{B}}}
\def\etC{{\etens{C}}}
\def\etD{{\etens{D}}}
\def\etE{{\etens{E}}}
\def\etF{{\etens{F}}}
\def\etG{{\etens{G}}}
\def\etH{{\etens{H}}}
\def\etI{{\etens{I}}}
\def\etJ{{\etens{J}}}
\def\etK{{\etens{K}}}
\def\etL{{\etens{L}}}
\def\etM{{\etens{M}}}
\def\etN{{\etens{N}}}
\def\etO{{\etens{O}}}
\def\etP{{\etens{P}}}
\def\etQ{{\etens{Q}}}
\def\etR{{\etens{R}}}
\def\etS{{\etens{S}}}
\def\etT{{\etens{T}}}
\def\etU{{\etens{U}}}
\def\etV{{\etens{V}}}
\def\etW{{\etens{W}}}
\def\etX{{\etens{X}}}
\def\etY{{\etens{Y}}}
\def\etZ{{\etens{Z}}}

% The true underlying data generating distribution
\newcommand{\pdata}{p_{\rm{data}}}
% The empirical distribution defined by the training set
\newcommand{\ptrain}{\hat{p}_{\rm{data}}}
\newcommand{\Ptrain}{\hat{P}_{\rm{data}}}
% The model distribution
\newcommand{\pmodel}{p_{\rm{model}}}
\newcommand{\Pmodel}{P_{\rm{model}}}
\newcommand{\ptildemodel}{\tilde{p}_{\rm{model}}}
% Stochastic autoencoder distributions
\newcommand{\pencode}{p_{\rm{encoder}}}
\newcommand{\pdecode}{p_{\rm{decoder}}}
\newcommand{\precons}{p_{\rm{reconstruct}}}

\newcommand{\laplace}{\mathrm{Laplace}} % Laplace distribution

\newcommand{\E}{\mathbb{E}}
\newcommand{\Ls}{\mathcal{L}}
\newcommand{\R}{\mathbb{R}}
\newcommand{\emp}{\tilde{p}}
\newcommand{\lr}{\alpha}
\newcommand{\reg}{\lambda}
\newcommand{\rect}{\mathrm{rectifier}}
\newcommand{\softmax}{\mathrm{softmax}}
\newcommand{\sigmoid}{\sigma}
\newcommand{\softplus}{\zeta}
\newcommand{\KL}{D_{\mathrm{KL}}}
\newcommand{\Var}{\mathrm{Var}}
\newcommand{\standarderror}{\mathrm{SE}}
\newcommand{\Cov}{\mathrm{Cov}}
% Wolfram Mathworld says $L^2$ is for function spaces and $\ell^2$ is for vectors
% But then they seem to use $L^2$ for vectors throughout the site, and so does
% wikipedia.
\newcommand{\normlzero}{L^0}
\newcommand{\normlone}{L^1}
\newcommand{\normltwo}{L^2}
\newcommand{\normlp}{L^p}
\newcommand{\normmax}{L^\infty}

\newcommand{\parents}{Pa} % See usage in notation.tex. Chosen to match Daphne's book.

\DeclareMathOperator*{\argmax}{arg\,max}
\DeclareMathOperator*{\argmin}{arg\,min}

\DeclareMathOperator{\sign}{sign}
\DeclareMathOperator{\Tr}{Tr}
\let\ab\allowbreak


\usepackage{amsmath}
\usepackage{amssymb}
\usepackage{amsthm}
\usepackage{mathtools}
\usepackage{hyperref}
\usepackage{url}
\usepackage[inkscapelatex=false]{svg}


\newcommand{\Prob}{\mathbb{P}}
\newcommand{\diag}{\operatorname{diag}}
\newcommand{\rank}{\operatorname{rank}}
\newcommand{\srank}{\operatorname{srank}}
\newcommand{\erank}{\operatorname{erank}}
\newcommand{\St}{\operatorname{St}}
\newcommand{\op}{\operatorname{op}}
\newcommand{\weak}{\Rightarrow}
\newcommand{\one}{\mathbf{1}}
\newcommand{\hadamard}{\mathbin{\odot}}
\newcommand{\bA}{\mathbf{A}}
\newcommand{\bB}{\mathbf{B}}
\newcommand{\bM}{\mathbf{M}}
\newcommand{\bL}{\mathbf{L}}
\newcommand{\bT}{\mathbf{T}}
\newcommand{\bQ}{\mathbf{Q}}
\newcommand{\bK}{\mathbf{K}}
\newcommand{\bU}{\mathbf{U}}
\newcommand{\bV}{\mathbf{V}}
\newcommand{\be}{\mathbf{e}}
\newcommand{\bu}{\mathbf{u}}
\newcommand{\bv}{\mathbf{v}}
\newcommand{\bw}{\mathbf{w}}
\newcommand{\bG}{\mathbf{G}}
\newcommand{\bH}{\mathbf{H}}
\newcommand{\bSigma}{\boldsymbol{\Sigma}}
\newcommand{\ESD}{\mu_{\mathrm{eig}}}
\newcommand{\SSD}{\mu_{\mathrm{sing}}}
\newcommand{\barPhi}{\overline{\Phi}}

\newtheorem{theorem}{Theorem}
\newtheorem{lemma}[theorem]{Lemma}
\newtheorem{proposition}[theorem]{Proposition}
\newtheorem{corollary}[theorem]{Corollary}
\newtheorem{conjecture}[theorem]{Conjecture}
\newtheorem{assumption}[theorem]{Assumption}
\theoremstyle{remark}
\newtheorem{remark}[theorem]{Remark}

\title{Spectral Analysis of Structured Causal Attention via State-Space Duality}

% Authors must not appear in the submitted version. They should be hidden
% as long as the \iclrfinalcopy macro remains commented out below.
\author{Mengda Li \\
       \And
       Zhaorui Dong \\
       \And
       Jianfeng Yao}

%\iclrfinalcopy

\begin{document}

\maketitle

\begin{abstract}
{
SSMs are increasingly viewed as replacements for, and competitors to,
Transformer models because of their linear inference cost.  A central question
is to understand their theoretical differences from Transformers, and we
approach this question through spectral analysis.  Following Mamba-2,
State-Space Duality allows us to unify SSMs and several related models within
the framework of structured causal attention.
}
\(
        \bM=\bL\hadamard \bQ\bK^\top,
\)
where $\bL$ is a structured causal mask.  
{
Self-attention behaves spectrally as an information compressor: the spikes in
the spectrum are preserved, while the bulk degenerates for long sequences.
By contrast, SSMs behave more like information diffusors.  Even low-rank
query-key matrices in trained networks can have diffuse spectral distributions
without degeneration.
}
\end{abstract}

\section{Introduction}
% {\color{blue}
% [SSMs are increasingly viewed as replacements for, and competitors to,
% Transformer models because of their linear inference cost.  A central question
% is to understand their theoretical differences from Transformers, and we
% approach this question through spectral analysis.]}
Transformers and state-space sequence models are now two central paradigms for
long-context modeling.  A useful way to compare them is not only by their
computational cost, but by how their token-token matrices move information
through a sequence.  Recent spectral work on softmax attention suggests an
``information compressor'' picture: after the softmax normalization, a small
number of spiked eigen/singular directions may survive, while the spectral bulk
degenerates \citep{li:2026spectral}.  This raises a natural
question for efficient alternatives to Transformers: do structured
state-space-style mixers compress information in the same way, or do their
spectra spread signal across many directions?

We study this question through the structured masked attention framework via State-Space Duality \citep{dao:2024transformers}.  Many recurrent or linear-time
sequence mixers can be written, at the level of their token-token interaction
matrix, as
\[
        \bM=\bL\hadamard \bQ\bK^\top ,
\]
where $\bQ\bK^\top$ is a query-key score matrix and $\bL$ is a structured
causal mask.  This representation does not mean that every such model is
literally an SSM.  Rather, it gives a common comparison class for SSMs,
retention models, and linear-attention baselines.  In this paper we focus on
three lower triangular masks: the causal all-one mask for linear attention
\citep{katharopoulos:2020transformers}, the one-semiseparable mask that appears
in SSD/Mamba-2 \citep{gu:2023mamba,dao:2024transformers}, and the Toeplitz
decay mask used in RetNet-style retention \citep{sun:2023retentive}.  General
Toeplitz and Fourier mixers are related structured alternatives, but the decay
Toeplitz mask is the simplest retention model for a first spectral comparison
\citep{qin:2023toeplitz,leethorp:2021fnet}.

Our guiding claim is that these structured causal masks are better understood
as information diffusers than as information compressors.  Even if the
pre-mask score matrix $\bQ\bK^\top$ is low rank, the Hadamard interaction with
$\bL$ can distribute this signal across the token-token matrix and prevent the
spectral law from collapsing to zero.  This is the theoretical contrast with
softmax attention that motivates the paper: softmax normalization tends to
concentrate spectral mass into a small number of surviving directions, whereas
structured causal masks can diffuse low-rank query-key signal into a
non-degenerate spectrum.

This paper begins with the low-rank regime for $\bQ$ and $\bK$.  This is a
trained or compressed model for the query-key scores, rather than a full
initialization model.  The assumption is deliberately asymmetric: the token-side
singular spaces may be dependent, as they are when $\bQ=XW_Q$ and $\bK=XW_K$,
but the right singular spaces in projection dimension are independent and
isotropic.  This isolates the loss created by extracting signal through random
right-singular overlaps.  The point is not to claim that low-rank scores alone
explain SSM behavior, but to show that the mask can turn low-rank scores into a
spectrally diffuse sequence mixer.

\subsection{Related Work}

The starting point is the Transformer attention mechanism of
\citet{vaswani:2017attention}.  Linear attention rewrites causal attention in a
form compatible with recurrent computation
\citep{katharopoulos:2020transformers}.  Selective state-space models and
Mamba introduced input-dependent SSM sequence mixers \citep{gu:2023mamba}, and
SSD/Mamba-2 showed that such models can be represented through a structured
state-space duality with semiseparable masked attention
\citep{dao:2024transformers}.  RetNet uses retention with a Toeplitz decay mask
\citep{sun:2023retentive}.  More general Toeplitz and Fourier token mixers
provide neighboring structured alternatives to dense attention
\citep{qin:2023toeplitz,leethorp:2021fnet}.

On the spectral side, recent random matrix analyses of softmax attention study
Gaussian equivalence, isotropic score geometry, and spectral degeneration
\citep{hayase:2025gaussian,li:2026spectral}.  Our work is motivated by the
complementary comparison: once softmax attention is viewed as an information
compressor, what spectral behavior is induced by the structured masks used in
SSM-style and retention models?  The Haar moment calculations below use the
orthogonal Weingarten calculus of \citet{collins:2006integration}.

\subsection{Contributions}

The paper makes three initial contributions.  First, we formulate the
compressor-versus-diffuser comparison in the common masked-score model
$\bM=\bL\hadamard\bQ\bK^\top$, with emphasis on causal linear attention,
one-semiseparable SSD/Mamba-2 masks, and Toeplitz decay masks for retention.
Second, in the low-rank model with independent Haar right singular spaces, we
prove an exact eigenvalue theorem for every lower triangular mask and identify
the common-Haar token threshold
$\sigma_a^Q\sigma_b^K\asymp \ell\sqrt d$.  Third, we show that the resulting
eigenvalue and singular-value behavior is non-degenerate: eigenvalues already
retain a macroscopic effective support through the diagonal overlaps, while
singular values expose the off-diagonal causal diffusion created by the mask.

\section{Generalized Masked Attention Model}

Let $\ell$ be the sequence length and $d$ the projection dimension.  We consider
\begin{equation}
    \bM=\bL\hadamard \bQ\bK^\top,\qquad
    \bQ,\bK\in\R^{\ell\times d},\quad
    \bL\in\R^{\ell\times \ell}.
    \label{eq:model}
\end{equation}
Here $\hadamard$ denotes the entrywise product.  The mask $\bL$ is the
structured object.  We focus on three unit-diagonal causal masks.

\begin{enumerate}
\item Causal linear attention.  The causal all-one mask is
\[
    L_{ij}=\one_{\{i\geq j\}} .
\]
Then $\bM$ is the unnormalized causal linear attention score matrix.

\item One-semiseparable SSM/SSD masks.  A scalar semiseparable causal mask is
\begin{equation}
    L_{ij}
    =
    \one_{\{i\geq j\}}\prod_{t=j+1}^{i}a_t,
    \qquad 0\leq a_t<1,
    \label{eq:semiseparable-mask}
\end{equation}
with the empty product equal to one.  In the SSD/Mamba-2 representation, this is
the structured attention form of a scalar state transition
$\mathbf A_t=a_t\mathbf I$, with $\bQ=\mathbf C$ and $\bK=\mathbf B$ in the
state-space notation.  Thus Mamba-type models appear here through the
semiseparable mask, not as a separate ambient model.

\item Toeplitz decay masks for retention.  A typical RetNet-type mask is
\[
    L_{ij}=\one_{\{i\geq j\}}\gamma^{i-j},
    \qquad 0<\gamma\leq 1,
\]
or a head-dependent version with $\gamma=\gamma_h$.  This is a stationary
Toeplitz mask and also a stationary one-semiseparable mask.
\end{enumerate}

Throughout the eigenvalue part of this paper, the main case is a lower
triangular mask:
\begin{equation}
    L_{ij}=0\quad\text{for }i<j.
    \label{eq:lower-triangular}
\end{equation}
The causal linear mask, the one-semiseparable SSM/SSD mask, and the Toeplitz
decay mask are all lower triangular and unit-diagonal.

\section{Low-Rank Model}

This section studies one regime of the general model: trained or compressed
query-key matrices whose effective ranks are small compared with the ambient
projection dimension.  Other regimes, including full-rank random initialization
and mask-specific singular-value laws, require separate analysis.

\subsection{Assumptions and Notation}

We write the singular value decompositions as
\begin{equation}
    \begin{aligned}
    \bQ&=\bU_Q\bSigma_Q\bV_Q^\top
        =\sum_{a=1}^{r_Q}\sigma^Q_a\bu^Q_a(\bv^Q_a)^\top,\\
    \bK&=\bU_K\bSigma_K\bV_K^\top
        =\sum_{b=1}^{r_K}\sigma^K_b\bu^K_b(\bv^K_b)^\top .
    \end{aligned}
    \label{eq:svd}
\end{equation}
where $\rank(\bQ)=r_Q$ and $\rank(\bK)=r_K$.  The columns of
$\bU_Q,\bU_K,\bV_Q,\bV_K$ are orthonormal, and
$\bSigma_Q,\bSigma_K$ contain the nonzero singular values.

\begin{assumption}[Independent Haar right spaces]
Given the token-side quantities and singular values, the matrices
\[
        \bV_Q\sim \mathrm{Haar}(\St(d,r_Q)),
        \qquad
        \bV_K\sim \mathrm{Haar}(\St(d,r_K))
\]
are independent.
\label{ass:haar}
\end{assumption}

\begin{assumption}[Common Haar token space]
Let $r\geq \max(r_Q,r_K)$.  The left singular vectors used by $\bQ$ and $\bK$
are columns of a common matrix
\[
    \bU=[\bu_1,\ldots,\bu_r]\sim \mathrm{Haar}(\St(\ell,r)),
\]
independent of the right singular spaces.  Thus
\[
    \bu^Q_a=\bu_a,\qquad
    \bu^K_b=\bu_b.
\]
\label{ass:common-left-haar}
\end{assumption}

Assumption~\ref{ass:haar} is the basic low-rank hypothesis: dependence can
remain on the token side, while the right singular spaces are independent in the
projection dimension.  Assumption~\ref{ass:common-left-haar} is stronger and is
used only when we want a calculable token-side null model.  It is motivated by
the fact that in attention one often has $\bQ=XW_Q$ and $\bK=XW_K$, so the left
singular spaces of $\bQ$ and $\bK$ are coupled through the same input $X$.
Here $\mathrm{Haar}(\St(n,r))$ denotes the uniform Stiefel law, equivalently the
law of the first $r$ columns of an $n\times n$ Haar orthogonal matrix.

\subsection{Eigenvalues of Lower Triangular Masks}

\begin{theorem}[Eigenvalues of lower triangular masked attention]
Assume $\bL$ is lower triangular and Assumptions~\ref{ass:haar} and
\ref{ass:common-left-haar} hold.  Suppose $r_Q,r_K$ are fixed while
$\ell,d\to\infty$, and
\[
    \theta_{ab}^{(d)}
    :=
    \frac{\sigma^Q_a\sigma^K_b}{\ell\sqrt d}
    \longrightarrow \theta_{ab}.
\]
For any fixed number of distinct tokens $i=1,\ldots,m$, jointly over those
tokens,
\begin{equation}
    \lambda_i(\bM)
    \weak
    L_{ii}\sum_{a=1}^{r_Q}\sum_{b=1}^{r_K}
    \theta_{ab}\,\xi_{i,a}\xi_{i,b}\,G_{ab},
    \qquad i=1,\ldots,m,
    \label{eq:common-haar-eigenvalue-limit}
\end{equation}
where the variables $\xi_{i,a}$ are independent $N(0,1)$ over token and
singular-vector indices, and $(G_{ab})$ is an independent matrix of independent
$N(0,1)$ entries.  The Gaussian matrix $(G_{ab})$ is shared across all tokens,
whereas the token-side variables $\xi_{i,a}$ depend on $i$.

In the causal linear attention, semiseparable SSM/SSD, and RetNet-type decay
models, $L_{ii}=1$.  Hence the limit reduces to
\[
    \lambda_i(\bM)
    \weak
    \sum_{a=1}^{r_Q}\sum_{b=1}^{r_K}
    \theta_{ab}\xi_{i,a}\xi_{i,b}G_{ab}.
\]
The natural nontrivial eigenvalue threshold in this unit-diagonal common-Haar
token model is $\sigma^Q_a\sigma^K_b\asymp \ell\sqrt d$.
\label{thm:lower-triangular-eigenvalues}
\end{theorem}

\begin{theorem}[Eigenvalue spectral phenomena]
Assume the setting of Theorem~\ref{thm:lower-triangular-eigenvalues},
$L_{ii}=1$, and $\boldsymbol{\theta}\neq0$.  Then the eigenvalues exhibit the
following count and rank phenomena.  Write
$\barPhi(t):=\Prob\{Z>t\}$ for $Z\sim N(0,1)$.

\begin{enumerate}
\item[(i)] Effective eigenvalue counts.  For every $\varepsilon>0$,
\begin{equation}
    \frac{\E\#\{i:|\lambda_i(\bM)|>\varepsilon\}}{\ell}
    \longrightarrow
    p_\varepsilon^{|\cdot|}(\boldsymbol{\theta})>0,
    \label{eq:expected-effective-count}
\end{equation}
where $p_\varepsilon^{|\cdot|}$ is defined in
\eqref{eq:eigenvalue-tail-definitions}.  The lower bound
\begin{equation}
    p_\varepsilon^{|\cdot|}(\boldsymbol{\theta})
    \geq
    \frac{3}{70}
    \barPhi\left(
    \frac{\sqrt2\,\varepsilon}{s_{\boldsymbol{\theta}}}
    \right)
    \label{eq:noncollapsed-esd-mass}
\end{equation}
follows from \eqref{eq:eigenvalue-tail-two-sided}, with
$s_{\boldsymbol{\theta}}^2=\sum_{a,b}
(1+2\mathbf 1_{\{a=b\}})\theta_{ab}^2$.  The Gaussian-chaos tail bound and
rank-one formula are given in \eqref{eq:eigenvalue-chaos-tail-bound} and
\eqref{eq:rank-one-tail-formula}.
Conditionally on the shared right-space Gaussian limit $\bG$, the empirical count
$\ell^{-1}\#\{i:|\lambda_i(\bM)|>\varepsilon\}$ converges to the corresponding
conditional Gaussian-chaos probability.

\item[(ii)] Eigenvalue stable rank and effective rank.  Define the eigenvalue
stable rank and eigenvalue effective rank by
\[
    \srank_{\lambda}(\bM)
    :=
    \frac{\sum_i|\lambda_i(\bM)|}{\max_i|\lambda_i(\bM)|},
    \qquad
    \erank_{\lambda}(\bM)
    :=
    \exp\left\{-\sum_i p_i\log p_i\right\},
\]
where
$p_i=|\lambda_i(\bM)|/\sum_j|\lambda_j(\bM)|$ and $0\log0=0$.  Then
\begin{equation}
    \srank_{\lambda}(\bM)
    =
    \Theta_{\Prob}\left(\frac{\ell}{\log\ell}\right),
    \qquad
    \erank_{\lambda}(\bM)=\Theta_{\Prob}(\ell).
    \label{eq:eigenvalue-rank-summary}
\end{equation}
Here $X_\ell=\Theta_{\Prob}(a_\ell)$ means that for every $\delta\in(0,1)$
there exist constants $0<c_\delta<C_\delta<\infty$ such that
$\Prob\{c_\delta a_\ell\leq X_\ell\leq C_\delta a_\ell\}\geq1-\delta$ for all
large $\ell$.  Sharper conditional prefactors and the entropy limit are stated
in \eqref{eq:eigenvalue-stable-rank-prefactor} and
\eqref{eq:eigenvalue-entropy-limit}.
\end{enumerate}
\label{thm:eigenvalue-spectral-phenomena}
\end{theorem}

\subsection{Singular Values of Unit-Diagonal Causal Masks}

Eigenvalues only see the diagonal of a lower triangular mask.  Singular values
see the full causal operator.  The following rank-one theorem records the
analogues of the three eigenvalue phenomena for causal linear attention, stable
one-semiseparable SSM/SSD masks, and stable decay masks.

\begin{theorem}[Singular-value spectral phenomena]
Assume $r_Q=r_K=1$, Assumptions~\ref{ass:haar},
\ref{ass:common-left-haar}, and suppose
\[
    \frac{\sigma^Q\sigma^K}{\ell\sqrt d}\longrightarrow \theta\neq0.
\]
Let $\bL$ be one of the following unit-diagonal causal masks: the causal linear
mask $L_{ij}=\one_{\{i\geq j\}}$, a decay mask
$L_{ij}=\one_{\{i\geq j\}}\gamma^{i-j}$ with $0\leq\gamma<1$, or a
one-semiseparable mask \eqref{eq:semiseparable-mask} with
$0\leq a_t\leq a<1$.  Let $C_{\mathrm{inv}}$ be a constant such that every
principal submatrix $\bL_S$ satisfies
$\|\bL_S^{-1}\|_{\op}\leq C_{\mathrm{inv}}$; for these masks such a constant
depends only on the mask family; see
Proposition~\ref{prop:principal-inverse-bound}.  Let
\[
    \mu_{\mathrm{sing},\ell}
    :=
    \frac{1}{\ell}\sum_{j=1}^{\ell}\delta_{s_j(\bM)}
\]
be the singular-value empirical measure, and define
\[
    \srank(\bM):=\frac{\|\bM\|_F^2}{\|\bM\|_{\op}^2},
    \qquad
    \erank_{\mathrm{sing}}(\bM)
    :=
    \exp\left\{-\sum_{j=1}^{\ell}q_j\log q_j\right\},
    \quad
    q_j=\frac{s_j(\bM)}{\sum_k s_k(\bM)}
\]
with $0\log0=0$.  Then the singular values exhibit the following phenomena.

\begin{enumerate}
\item[(i)] Non-collapse of the singular-value empirical law.  Let
$G\sim N(0,1)$ denote the right-space limit
$\sqrt d\,\langle\bv^Q,\bv^K\rangle\weak G$.  Conditionally on the event
$G\neq0$, $\mu_{\mathrm{sing},\ell}$ does not converge to $\delta_0$.  More
precisely, for every $t>0$,
\begin{equation}
    \liminf_{\ell\to\infty}
    \mu_{\mathrm{sing},\ell}
    \left(
        \left(
        \frac{|\theta G|t^2}{2C_{\mathrm{inv}}},
        \infty
        \right)
    \right)
    \geq
    \Prob\{|N(0,1)|>t\}
    \label{eq:singular-noncollapse}
\end{equation}
in probability.

\item[(ii)] Stable rank and effective rank for stable masks.  For the decay and
stable one-semiseparable masks,
\begin{equation}
    \srank(\bM)
    \geq
    c\,\frac{\ell}{(\log\ell)^2},
    \qquad
    \erank_{\mathrm{sing}}(\bM)
    \geq
    c\,\frac{\ell}{\log\ell}
    \label{eq:stable-mask-singular-ranks}
\end{equation}
with probability tending to one, for a constant $c>0$ depending only on the
mask parameters.

\item[(iii)] Stable rank and effective rank for causal linear attention.  For
the causal linear mask, assume the weighted Volterra spread estimate
\eqref{eq:weighted-volterra-spread}.  Then
\[
    \srank(\bM)=\Theta_{\Prob}(1).
\]
Moreover,
\begin{equation}
    \erank_{\mathrm{sing}}(\bM)\longrightarrow\infty
    \label{eq:linear-mask-singular-ranks}
\end{equation}
in probability; in particular
$\erank_{\mathrm{sing}}(\bM)\geq c\log\ell$ with probability tending to one.
\end{enumerate}
\label{thm:singular-spectral-phenomena}
\end{theorem}

\section{Experimentation}

\begin{figure}[t]
  \centering
  \includesvg[width=\linewidth]{figure/layer0_m_rank_vs_length_iqr2.svg}
  \caption{Stable rank and effective rank of the layer-0 structured attention
  matrices from the Mamba-2 
  \texttt{mistralai/Mamba-Codestral-7B-v0.1} on Hugging Face, evaluated on the
  Wikitext-103 validation split as a function of sequence length.  Curves show
  median trends; shaded bands indicate interquartile ranges and wider
  percentile intervals.}
  \label{fig:rank-wikitext}
\end{figure}

\subsubsection*{Acknowledgments}
Funding and competing-interest disclosures to be added.

\bibliography{ref}
\bibliographystyle{plainnat}

\appendix

\section{Proofs for the Low-Rank Model}

\begin{lemma}[Haar overlap scale]
Under Assumption~\ref{ass:haar}, for deterministic
$\mathbf a\in\R^{r_Q}$ and $\mathbf b\in\R^{r_K}$,
\[
    \mathbf a^\top\bH\mathbf b
    \stackrel{d}{=}
    \|\mathbf a\|\,\|\mathbf b\|\,\xi_d,
\]
where $\xi_d$ is the inner product of two independent uniform unit vectors in
$\R^d$.  In particular,
\[
    \E\,\mathbf a^\top\bH\mathbf b=0,
    \qquad
    \E|\mathbf a^\top\bH\mathbf b|^2
    =
    \frac{\|\mathbf a\|^2\|\mathbf b\|^2}{d}.
\]
Moreover, for a universal constant $c>0$,
\[
    \Prob\left(
    |\mathbf a^\top\bH\mathbf b|
    >
    t\|\mathbf a\|\|\mathbf b\|
    \right)
    \leq 2\exp(-cdt^2),
    \qquad 0\leq t\leq 1.
\]
If $r_Q,r_K$ are fixed and $d\to\infty$, then
\[
        \sqrt d\,\bH\weak \bG,
\]
where $\bG\in\R^{r_Q\times r_K}$ has independent $N(0,1)$ entries.
\label{lem:haar-overlap}
\end{lemma}

\begin{proof}
The vectors $\bV_Q\mathbf a$ and $\bV_K\mathbf b$ are independent uniform
vectors on spheres of radii $\|\mathbf a\|$ and $\|\mathbf b\|$ in $\R^d$.
Therefore
\[
    \mathbf a^\top\bH\mathbf b
    =
    \langle \bV_Q\mathbf a,\bV_K\mathbf b\rangle
\]
has the stated distribution.  The variance and concentration bound are standard
facts for a coordinate of a uniform vector on the sphere.  The finite-dimensional
Gaussian limit follows from the Gaussian representation of Haar frames; for
fixed ranks, the normalization and orthogonalization errors vanish as
$d\to\infty$.
\end{proof}

\begin{lemma}[Token-side Haar products]
Under Assumption~\ref{ass:common-left-haar}, for any fixed token $i$ and fixed
indices $a,b$,
\[
    \E\,u_a(i)u_b(i)=\delta_{ab}\,\mathrm{Wg}^{O(\ell)}(1)
    =
    \frac{\delta_{ab}}{\ell}.
\]
Moreover,
\[
    \E\,u_a(i)^2u_b(i)^2
    =
    \begin{cases}
    \displaystyle \frac{3}{\ell(\ell+2)}, & a=b,\\[1ex]
    \displaystyle \frac{1}{\ell(\ell+2)}, & a\neq b.
    \end{cases}
\]
Equivalently, for fixed $m,r$ and distinct fixed tokens $i=1,\ldots,m$,
\[
    \left(\sqrt{\ell}\,(u_1(i),\ldots,u_r(i))\right)_{i=1}^m
    \weak
    \left((\xi_{i,1},\ldots,\xi_{i,r})\right)_{i=1}^m,
    \qquad
    \xi_{i,a}\stackrel{\mathrm{i.i.d.}}{\sim}N(0,1),
\]
and hence
\[
    \left(\ell\,u_a(i)u_b(i)\right)_{i,a,b}
    \weak
    \left(\xi_{i,a}\xi_{i,b}\right)_{i,a,b}.
\]
\label{lem:token-haar-products}
\end{lemma}

\begin{proof}
Let $O\in O(\ell)$ be Haar.  The first $r$ columns of $O$ have the same law as
$\bU$.  We apply the orthogonal Weingarten formula of
\citet{collins:2006integration}.  For two entries,
\[
    \E O_{ia}O_{ib}
    =
    \delta_{ab}\,\mathrm{Wg}^{O(\ell)}(1)
    =
    \frac{\delta_{ab}}{\ell}.
\]
For four entries,
\[
    \E O_{ia}O_{ia}O_{ib}O_{ib}
    =
    \sum_{\pi,\sigma\in\mathcal P_2(4)}
    \delta_i^\pi\,\delta_{aabb}^{\sigma}\,
    \mathrm{Wg}^{O(\ell)}(\pi^{-1}\sigma),
\]
where $\mathcal P_2(4)$ is the set of pair partitions of four points.  Since all
row indices are equal, $\delta_i^\pi=1$ for every $\pi$.  Evaluating the three
pairings gives
\[
    \E O_{ia}^2O_{ib}^2
    =
    \frac{1+2\delta_{ab}}{\ell(\ell+2)}.
\]
This is exactly the displayed formula.  The finite-row Gaussian limit follows
from the standard asymptotic form of the same Weingarten calculus, equivalently
from the Gaussian representation of finitely many rows of a Haar orthogonal
matrix.
\end{proof}

\begin{lemma}[Empirical Haar-row law of large numbers]
Let $r$ be fixed and let
$\bU_\ell=[\bu_1,\ldots,\bu_r]\sim\mathrm{Haar}(\St(\ell,r))$.  Define
\[
    \mathbf z_i
    :=
    \sqrt\ell\,(u_1(i),\ldots,u_r(i))\in\R^r,
    \qquad
    \mu_\ell
    :=
    \frac1\ell\sum_{i=1}^{\ell}\delta_{\mathbf z_i}.
\]
Then $\mu_\ell\weak N(0,\mathbf I_r)$ in probability.  Equivalently, for every
continuous $f:\R^r\to\R$ of polynomial growth,
\[
    \frac1\ell\sum_{i=1}^{\ell}f(\mathbf z_i)
    \longrightarrow
    \E f(\boldsymbol{\xi}),
    \qquad
    \boldsymbol{\xi}\sim N(0,\mathbf I_r),
\]
in probability.
\label{lem:empirical-haar-row}
\end{lemma}

\begin{proof}
Let $\bG_\ell\in\R^{\ell\times r}$ have independent $N(0,1)$ entries, with rows
$\mathbf g_i$.  The Gaussian representation of the Haar Stiefel matrix gives
\[
    \bU_\ell
    \stackrel{d}{=}
    \bG_\ell(\bG_\ell^\top\bG_\ell)^{-1/2}.
\]
Hence
\[
    \mathbf z_i
    \stackrel{d}{=}
    \mathbf g_i
    \left(\frac{\bG_\ell^\top\bG_\ell}{\ell}\right)^{-1/2}.
\]
Since $r$ is fixed,
$\ell^{-1}\bG_\ell^\top\bG_\ell\to\mathbf I_r$ almost surely.  The empirical
law of the independent rows $\mathbf g_i$ converges weakly to
$N(0,\mathbf I_r)$ almost surely, and the common multiplier above converges to
the identity.  This proves weak convergence of $\mu_\ell$.  For functions of
polynomial growth, truncate $f$ outside a large ball and use the convergence of
Gaussian moments of fixed order to remove the truncation.
\end{proof}

\paragraph{Exact diagonal formulas.}
For each token $i$, define the left singular loading vectors
\begin{equation}
    \bar{\bu}_{Q,i}
    :=
    \bSigma_Q\bU_Q^\top\be_i\in\R^{r_Q},
    \qquad
    \bar{\bu}_{K,i}
    :=
    \bSigma_K\bU_K^\top\be_i\in\R^{r_K}.
    \label{eq:bar-loadings}
\end{equation}
These vectors include both singular values and token-side leverage.  Let
\begin{equation}
    \bH:=\bV_Q^\top\bV_K\in\R^{r_Q\times r_K}.
    \label{eq:H}
\end{equation}
Then every entry of $\bQ\bK^\top$ has the form
\begin{equation}
    (\bQ\bK^\top)_{ij}
    =
    \bar{\bu}_{Q,i}^{\top}\bH\,\bar{\bu}_{K,j}.
    \label{eq:entry-form}
\end{equation}

\begin{proposition}[Exact eigenvalues for lower triangular masks]
Assume $\bL$ is lower triangular.  Then the eigenvalues of
$\bM=\bL\hadamard\bQ\bK^\top$ are exactly
\begin{equation}
    \lambda_i(\bM)
    =
    L_{ii}\bar{\bu}_{Q,i}^{\top}\bH\,\bar{\bu}_{K,i},
    \qquad i=1,\ldots,\ell.
    \label{eq:eigenvalues-exact}
\end{equation}
Equivalently,
\begin{equation}
    \lambda_i(\bM)
    =
    L_{ii}
    \sum_{a=1}^{r_Q}\sum_{b=1}^{r_K}
    \sigma^Q_a\sigma^K_b\,
    u^Q_a(i)u^K_b(i)\,
    \langle\bv^Q_a,\bv^K_b\rangle .
    \label{eq:eigenvalue-expansion}
\end{equation}
If Assumption~\ref{ass:haar} also holds and $r_Q,r_K$ are fixed as
$d\to\infty$, then with
\[
    Z_{ab}^{(d)}
    :=
    \sqrt d\,\langle\bv^Q_a,\bv^K_b\rangle,
    \qquad
    c_{ab}^{(d)}(i)
    :=
    \frac{L_{ii}\sigma^Q_a\sigma^K_bu^Q_a(i)u^K_b(i)}{\sqrt d},
\]
one has the exact representation
\begin{equation}
    \lambda_i(\bM)
    =
    \sum_{a=1}^{r_Q}\sum_{b=1}^{r_K}
    c_{ab}^{(d)}(i)Z_{ab}^{(d)}.
    \label{eq:fixed-rank-exact-characterization}
\end{equation}
Moreover,
\[
    (Z_{ab}^{(d)})_{a,b}\weak (G_{ab})_{a,b},
\]
where the entries $G_{ab}$ are independent $N(0,1)$.
\label{prop:exact-lower-triangular-eigenvalues}
\end{proposition}

\begin{proof}
Since $\bL$ is lower triangular, $\bM=\bL\hadamard\bQ\bK^\top$ is lower
triangular.  The eigenvalues of a triangular matrix are its diagonal entries,
which gives \eqref{eq:eigenvalues-exact}.  Expanding \eqref{eq:svd} gives
\eqref{eq:eigenvalue-expansion}.  The exact representation
\eqref{eq:fixed-rank-exact-characterization} is a rewriting of
\eqref{eq:eigenvalue-expansion}, and the convergence of
$(Z_{ab}^{(d)})_{a,b}$ follows from Lemma~\ref{lem:haar-overlap}.
\end{proof}

\begin{proof}[Proof of Theorem~\ref{thm:lower-triangular-eigenvalues}]
By Proposition~\ref{prop:exact-lower-triangular-eigenvalues} and
Assumption~\ref{ass:common-left-haar},
\[
    \lambda_i(\bM)
    =
    L_{ii}\sum_{a=1}^{r_Q}\sum_{b=1}^{r_K}
    \frac{\sigma^Q_a\sigma^K_b}{\ell\sqrt d}
    \{\ell u_a(i)u_b(i)\}
    Z_{ab}^{(d)}.
\]
Lemma~\ref{lem:token-haar-products} gives
\[
\left(\ell u_a(i)u_b(i)\right)_{i,a,b}
\weak
\left(\xi_{i,a}\xi_{i,b}\right)_{i,a,b},
\]
and Lemma~\ref{lem:haar-overlap} gives
$Z_{ab}^{(d)}\weak G_{ab}$ jointly over fixed ranks.  The two limits are
independent by assumption, so Slutsky's theorem proves
\eqref{eq:common-haar-eigenvalue-limit}.  The unit-diagonal special case follows
by setting $L_{ii}=1$ in this limit.
\end{proof}

\subsection{Eigenvalue Details}
\label{app:eigenvalue-details}

\begin{proposition}[Non-collapse of the eigenvalue empirical law]
Assume the setting of Theorem~\ref{thm:lower-triangular-eigenvalues},
$L_{ii}=1$, and $\boldsymbol{\theta}\neq0$.  Let
$\mu_{\mathrm{eig},\ell}:=\ell^{-1}\sum_{i=1}^{\ell}
\delta_{\lambda_i(\bM)}$.  Conditionally on the shared right-space Gaussian
limit $\bG$,
\begin{equation}
    \mu_{\mathrm{eig},\ell}
    \weak
    \nu_{\bG}
    :=
    \mathcal L_{\boldsymbol{\xi}}
    \left(
    \boldsymbol{\xi}^\top\bB(\bG)\boldsymbol{\xi}
    \right),
    \label{eq:noncollapsed-esd-limit}
\end{equation}
where $\bB(\bG)$ is the symmetric part of the matrix
$A_{ab}(\bG)=\theta_{ab}G_{ab}$.  Moreover
$\nu_{\bG}\neq\delta_0$ almost surely, so
$\mu_{\mathrm{eig},\ell}$ does not converge to $\delta_0$.
\label{prop:eigenvalue-esd-noncollapse}
\end{proposition}

\begin{proof}
Condition on $\bG$.  Lemma~\ref{lem:empirical-haar-row} gives empirical
convergence of the diagonal eigenvalues to the law of
$\boldsymbol{\xi}^\top\bB(\bG)\boldsymbol{\xi}$, which proves
\eqref{eq:noncollapsed-esd-limit}.  If $\boldsymbol{\theta}\neq0$, then
$\bB(\bG)\neq0$ almost surely, so the quadratic form is not identically zero and
$\nu_{\bG}\neq\delta_0$.
\end{proof}

\begin{proof}[Proof of Theorem~\ref{thm:eigenvalue-spectral-phenomena}]
Theorem~\ref{thm:lower-triangular-eigenvalues} gives, for a fixed token,
\[
    \lambda_i(\bM)
    \weak
    \sum_{a=1}^{r_Q}\sum_{b=1}^{r_K}
    \theta_{ab}\xi_a\xi_bG_{ab},
\]
where $(G_{ab})$ is independent of $\boldsymbol{\xi}$.  Conditional on
$\boldsymbol{\xi}$, the right-hand side is a centered Gaussian random variable
with variance
\[
    V_{\boldsymbol{\theta}}(\boldsymbol{\xi})
    =
    \sum_{a=1}^{r_Q}\sum_{b=1}^{r_K}
    \theta_{ab}^2\xi_a^2\xi_b^2.
\]
Set
\[
    s_{\boldsymbol{\theta}}^2
    =
    \E V_{\boldsymbol{\theta}}(\boldsymbol{\xi})
    =
    \sum_{a=1}^{r_Q}\sum_{b=1}^{r_K}
    (1+2\mathbf 1_{\{a=b\}})\theta_{ab}^2.
\]
For $\varepsilon>0$, define the one-sided and two-sided limiting tail
probabilities by
\begin{equation}
    p_\varepsilon^+(\boldsymbol{\theta})
    =
    \E_{\boldsymbol{\xi}}
    \barPhi\left(
        \frac{\varepsilon}
        {\sqrt{V_{\boldsymbol{\theta}}(\boldsymbol{\xi})}}
    \right),
    \qquad
    p_\varepsilon^{|\cdot|}(\boldsymbol{\theta})
    =
    2p_\varepsilon^+(\boldsymbol{\theta}),
    \label{eq:eigenvalue-tail-definitions}
\end{equation}
with the integrand set equal to zero when
$V_{\boldsymbol{\theta}}(\boldsymbol{\xi})=0$.
Therefore
\[
    \Prob_{\bG}
    \left\{
    \sum_{a,b}\theta_{ab}\xi_a\xi_bG_{ab}>\varepsilon
    \,\middle|\,\boldsymbol{\xi}
    \right\}
    =
    \barPhi\left(
        \frac{\varepsilon}{\sqrt{V_{\boldsymbol{\theta}}(\boldsymbol{\xi})}}
    \right),
\]
with value zero when the variance is zero.  Averaging over
$\boldsymbol{\xi}$ gives the one-sided formula defining
$p_\varepsilon^+(\boldsymbol{\theta})$.  Since the limiting variable is
symmetric about zero after averaging over $\bG$, the two-sided formula is twice
the one-sided formula.

The truncation estimate follows by splitting on
$V_{\boldsymbol{\theta}}(\boldsymbol{\xi})\leq u^2$.  On this event the
integrand is at most $\barPhi(\varepsilon/u)$, while the complement has
probability at most
$s_{\boldsymbol{\theta}}^2/u^2$ by Markov's inequality.  Therefore, for every
$u>0$,
\begin{equation}
    \frac{3}{70}
    \barPhi\left(
    \frac{\sqrt2\,\varepsilon}{s_{\boldsymbol{\theta}}}
    \right)
    \leq
    p_\varepsilon^{|\cdot|}(\boldsymbol{\theta})
    \leq
    2\barPhi(\varepsilon/u)+2s_{\boldsymbol{\theta}}^2/u^2.
    \label{eq:eigenvalue-tail-two-sided}
\end{equation}
The Gaussian chaos bound is the standard hypercontractive tail estimate for a
homogeneous degree-three Gaussian chaos, applied to
$\sum_{a,b}\theta_{ab}\xi_a\xi_bG_{ab}$:
\begin{equation}
    p_\varepsilon^{|\cdot|}(\boldsymbol{\theta})
    \leq
    C\exp\left\{
        -c\left(\frac{\varepsilon}{s_{\boldsymbol{\theta}}}\right)^{2/3}
    \right\},
    \label{eq:eigenvalue-chaos-tail-bound}
\end{equation}
for universal constants $c,C>0$.

For the lower estimate, Paley--Zygmund applied to the nonnegative random
variable $V_{\boldsymbol{\theta}}(\boldsymbol{\xi})$ gives
\[
    \Prob\left\{
    V_{\boldsymbol{\theta}}(\boldsymbol{\xi})
    \geq
    \frac{s_{\boldsymbol{\theta}}^2}{2}
    \right\}
    \geq
    \frac{1}{4K_{\boldsymbol{\theta}}}.
\]
On this event,
$\barPhi\{\varepsilon/V_{\boldsymbol{\theta}}(\boldsymbol{\xi})^{1/2}\}
\geq
\barPhi(\sqrt2\,\varepsilon/s_{\boldsymbol{\theta}})$, which proves
the lower bound in \eqref{eq:noncollapsed-esd-mass}.  Finally,
$K_{\boldsymbol{\theta}}\leq 35/3$ because
\[
    \frac{
    \E\{\xi_a^2\xi_b^2\xi_{a'}^2\xi_{b'}^2\}}
    {(1+2\mathbf 1_{\{a=b\}})(1+2\mathbf 1_{\{a'=b'\}})}
    \leq
    \frac{35}{3}
\]
for all indices; the maximum occurs when $a=b=a'=b'$.

Linearity of expectation gives
\[
    \E N_\varepsilon^+
    =
    \sum_{i=1}^{\ell}\Prob\{\lambda_i(\bM)>\varepsilon\},
\]
and similarly for $N_\varepsilon^{|\cdot|}$.  The limit in
\eqref{eq:expected-effective-count} follows from the fixed-token convergence
and exchangeability of the Haar token rows.  Finally, if one keeps the shared
right-space limit $\bG$ fixed, Lemma~\ref{lem:empirical-haar-row} applied to
the token variables gives convergence of the empirical fractions to the
corresponding conditional Gaussian-chaos tail probabilities.  The indicator
functions are handled by approximation at continuity points; for
$\varepsilon>0$ the relevant boundaries have probability zero unless the
limiting quadratic form is identically zero, which has probability zero when
$\boldsymbol{\theta}\neq0$.

In the rank-one case,
\begin{equation}
    p_\varepsilon^+(\theta)
    =
    \E_{\xi}
    \barPhi\left(\frac{\varepsilon}{|\theta|\xi^2}\right)
    =
    \int_0^\infty
    \barPhi\left(\frac{\varepsilon}{|\theta|x}\right)
    \frac{e^{-x/2}}{\sqrt{2\pi x}}\,dx .
    \label{eq:rank-one-tail-formula}
\end{equation}
Writing $x=\xi^2$ and using the $\chi^2(1)$ density gives the rank-one
one-dimensional formula.

Condition on $\bG$.  Define $\bA(\bG)$ by
$A_{ab}(\bG)=\theta_{ab}G_{ab}$ and let
$\bB(\bG)=\{\bA(\bG)+\bA(\bG)^\top\}/2$.  Write
\[
    Y_{\bG}=\boldsymbol{\xi}^\top\bB(\bG)\boldsymbol{\xi},
\]
because the limiting eigenvalue quadratic form only sees the symmetric part of
$\bA(\bG)$.  Lemma~\ref{lem:empirical-haar-row} gives empirical convergence of
the token eigenvalues to independent copies of $Y_{\bG}$.

Diagonalizing $\bB(\bG)$ gives
\[
    \boldsymbol{\xi}^\top\bB(\bG)\boldsymbol{\xi}
    \stackrel{d}{=}
    \sum_{k=1}^r\rho_k(\bG)\chi_k^2.
\]
The largest absolute value among $\ell$ independent copies of this finite-rank
Gaussian quadratic form has scale $2\rho_*(\bG)\log\ell$ whenever
$\rho_*(\bG)>0$.  Also,
\[
    \frac{1}{\ell}\sum_{i=1}^{\ell}|\lambda_i(\bM)|
    \weak
    \mu_1(\bG).
\]
Combining these two facts gives
\begin{equation}
    \frac{\log\ell}{\ell}r_{\lambda,1}(\bM)
    \weak
    \frac{\mu_1(\bG)}{2\rho_*(\bG)}.
    \label{eq:eigenvalue-stable-rank-prefactor}
\end{equation}
The upper bound on the prefactor follows from
\[
    \mu_1(\bG)
    \leq
    \E_{\boldsymbol{\xi}}\sum_{k=1}^r|\rho_k(\bG)|\chi_k^2
    =
    \sum_{k=1}^r|\rho_k(\bG)|
    \leq
    r\rho_*(\bG).
\]
This gives the displayed upper bound after dividing by $2\rho_*(\bG)$.  For
the lower bound, the Paley--Zygmund estimate used above and the fourth-moment
comparison for Gaussian quadratic forms give the universal bound
$\mu_1(\bG)\geq \sqrt2\,\rho_*(\bG)/9$ whenever $\bB(\bG)\neq0$.  This proves
the $\Theta_{\Prob}(\ell/\log\ell)$ statement because $\bB(\bG)\neq0$ almost
surely when $\boldsymbol{\theta}\neq0$.

For the entropy statement, condition on $\bG$.  The token variables converge to
independent copies of $Y_{\bG}=\boldsymbol{\xi}^\top\bB(\bG)\boldsymbol{\xi}$.
The law of large numbers gives
\[
    \frac{1}{\ell}\sum_{i=1}^{\ell}|\lambda_i(\bM)|
    \weak
    \mu_1(\bG),
\]
and
\[
    \frac{1}{\ell}\sum_{i=1}^{\ell}
    |\lambda_i(\bM)|
    \log |\lambda_i(\bM)|
    \weak
    \eta_1(\bG).
\]
Since
\[
    -\sum_{i=1}^{\ell}p_i\log p_i
    =
    \log\left(\sum_{i=1}^{\ell}w_i\right)
    -
    \frac{\sum_{i=1}^{\ell}w_i\log w_i}
    {\sum_{i=1}^{\ell}w_i},
\]
we obtain
\begin{equation}
    \log\operatorname{erank}_{\lambda}(\bM)
    =
    \log\ell+\log\mu_1(\bG)
    -
    \frac{\eta_1(\bG)}{\mu_1(\bG)}
    +o_{\Prob}(1),
    \label{eq:eigenvalue-entropy-limit}
\end{equation}
which gives the entropy-rank limit.  Finally, the Gaussian-chaos entropy
estimate gives
\[
    \log\mu_1(\bG)
    -
    \frac{\eta_1(\bG)}{\mu_1(\bG)}
    \geq
    -\log81,
\]
and the entropy of a probability vector on $\ell$ points is at most $\log\ell$.
This proves the linear lower and upper bounds in
\eqref{eq:eigenvalue-rank-summary}.
\end{proof}

\begin{remark}[Signal extraction threshold comparison]
For the structured causal masks above, the common-Haar token calculation gives
the eigenvalue-level requirement
$\sigma^Q_a\sigma^K_b \gg \ell\sqrt d$.  This is much stronger than the
softmax-attention threshold in \citet{li:2026spectral}.  In that setting,
isotropic score geometry can preserve signals of order larger than $\log \ell$.
\label{rem:threshold-comparison}
\end{remark}

\begin{remark}[Why $\rho_*$ is not one]
The quantity $\rho_*(\bG)$ above is not equal to one in this fixed-rank setting.
It is the largest absolute eigenvalue of the finite symmetric variance-profile
matrix $\bB(\bG)=\{\bA(\bG)+\bA(\bG)^\top\}/2$, where
$A_{ab}(\bG)=\theta_{ab}G_{ab}$.  The circular law concerns the eigenvalues of
large non-Hermitian Ginibre matrices after a $1/\sqrt r$ normalization.  It does
not apply here.
\label{rem:rho-star-not-one}
\end{remark}

\subsection{Singular-Value Details}
\label{app:singular-details}

\begin{proposition}[Low-rank singular-value reduction]
Under the decomposition \eqref{eq:svd},
\begin{equation}
    \bM
    =
    \sum_{a=1}^{r_Q}
    \sum_{b=1}^{r_K}
    \sigma^Q_a\sigma^K_b
    \langle \bv^Q_a,\bv^K_b\rangle
    \diag(\bu^Q_a)\,\bL\,\diag(\bu^K_b).
    \label{eq:singular-reduction}
\end{equation}
In the rank-one case,
\[
    s_j(\bM)
    =
    |\sigma^Q\sigma^K\langle \bv^Q,\bv^K\rangle|\,
    s_j\{\diag(\bu^Q)\bL\diag(\bu^K)\}.
\]
\label{prop:singular-reduction}
\end{proposition}

\begin{proof}
The identity
\[
    \bQ\bK^\top
    =
    \sum_{a=1}^{r_Q}
    \sum_{b=1}^{r_K}
    \sigma^Q_a\sigma^K_b
    \langle \bv^Q_a,\bv^K_b\rangle
    \bu^Q_a(\bu^K_b)^\top
\]
combined with
\[
    \bL\hadamard(\bu^Q_a(\bu^K_b)^\top)
    =
    \diag(\bu^Q_a)\bL\diag(\bu^K_b)
\]
gives the claim.
\end{proof}

\begin{proposition}[Stable semiseparable norm bound]
Let $\bL$ be the semiseparable mask \eqref{eq:semiseparable-mask} and assume
$0\leq a_t\leq a<1$ for all $t$.  Then
$\|\bL\|_{\op}\leq(1-a)^{-1}$.
\label{prop:semiseparable-bound}
\end{proposition}

\begin{proof}
For $i\geq j$, $|L_{ij}|\leq a^{i-j}$.  Hence every row sum and every column sum
of $|\bL|$ is bounded by $\sum_{k\geq0}a^k=(1-a)^{-1}$.  Schur's test gives
$\|\bL\|_{\op}\leq(1-a)^{-1}$.
\end{proof}

\begin{proposition}[Principal inverse bounds for causal masks]
Let $\bL_S$ denote the principal submatrix of $\bL$ indexed by
$S\subset\{1,\ldots,\ell\}$, with the inherited order.  For the causal linear
mask, $\|\bL_S^{-1}\|_{\op}\leq2$.  For the decay mask
$L_{ij}=\one_{\{i\geq j\}}\gamma^{i-j}$ with $0\leq\gamma<1$,
$\|\bL_S^{-1}\|_{\op}\leq1+\gamma$.  For the semiseparable mask
\eqref{eq:semiseparable-mask} with $0\leq a_t\leq a<1$,
$\|\bL_S^{-1}\|_{\op}\leq1+a$.
\label{prop:principal-inverse-bound}
\end{proposition}

\begin{proof}
Write $S=\{i_1<\cdots<i_m\}$.  In the causal linear case, $\bL_S$ is the
$m\times m$ lower triangular all-one matrix.  Its inverse has one on the
diagonal, minus one on the first subdiagonal, and zero elsewhere, so its
operator norm is at most two.

For the decay mask, $(\bL_S)_{pq}=\one_{\{p\geq q\}}\gamma^{i_p-i_q}$.  Its
inverse is lower bidiagonal with diagonal entries one and subdiagonal entries
$-\gamma^{i_p-i_{p-1}}$, whose absolute values are at most $\gamma$.  The
semiseparable case is identical after replacing
$\gamma^{i_p-i_{p-1}}$ by $\prod_{t=i_{p-1}+1}^{i_p}a_t$, whose absolute value
is at most $a$.  The displayed operator norm bounds follow from the triangle
inequality.
\end{proof}

\begin{assumption}[Weighted Volterra spread estimate]
For the causal linear mask $\bL_{ij}=\one_{\{i\geq j\}}$, write the rank-one
token vector as $\bu=(u(1),\ldots,u(\ell))$ and
$\bw=\sqrt\ell(|u(1)|,\ldots,|u(\ell)|)$.  The causal-linear part of
Theorem~\ref{thm:singular-spectral-phenomena} assumes that, with probability
tending to one, there exist constants $c_V,C_V>0$ such that
\begin{equation}
    \|\diag(\bw)\bL\diag(\bw)\|_*
    \geq c_V\ell\log\ell,
    \qquad
    \|\diag(\bw)\bL\diag(\bw)\|_{\op}
    \leq C_V\ell .
    \label{eq:weighted-volterra-spread}
\end{equation}
\label{ass:weighted-volterra-spread}
\end{assumption}

\begin{proof}[Proof of Theorem~\ref{thm:singular-spectral-phenomena}]
By Proposition~\ref{prop:singular-reduction}, in the rank-one case,
\[
    s_j(\bM)
    =
    |\sigma^Q\sigma^K\langle\bv^Q,\bv^K\rangle|\,
    s_j\{\diag(\bu)\bL\diag(\bu)\}.
\]
Let $Z_d=\sqrt d\,\langle\bv^Q,\bv^K\rangle$.  Assumption~\ref{ass:haar} gives
$Z_d\weak G$, and the signal scaling gives
\[
    \frac{\sigma^Q\sigma^K\langle\bv^Q,\bv^K\rangle}{\ell}
    \weak
    \theta G.
\]
Writing $w_i=\sqrt\ell\,|u(i)|$ and using sign matrices, the singular values of
$\diag(\bu)\bL\diag(\bu)$ are $\ell^{-1}$ times those of
$\diag(\bw)\bL\diag(\bw)$.  Hence the singular-value problem reduces to the
weighted mask
\[
    \bT_\ell:=\diag(\bw)\bL\diag(\bw),
    \qquad
    w_i=\sqrt\ell\,|u(i)|.
\]

We first prove non-collapse.  Fix $t>0$ and set
$S_t=\{i:w_i>t\}$.  Lemma~\ref{lem:empirical-haar-row}, applied to the
continuity set $\{\mathbf z:|z_1|>t\}$, gives
$|S_t|/\ell\to\Prob\{|N(0,1)|>t\}$.  The principal submatrix
$(\bT_\ell)_{S_t}$ satisfies
\[
    s_{\min}\{(\bT_\ell)_{S_t}\}
    \geq
    t^2s_{\min}(\bL_{S_t})
    \geq
    \frac{t^2}{C_{\mathrm{inv}}}.
\]
By singular-value interlacing for submatrices, $\bT_\ell$ has at least
$|S_t|$ singular values larger than $t^2/C_{\mathrm{inv}}$.  Since
$|\sigma^Q\sigma^K\langle\bv^Q,\bv^K\rangle|/\ell\to|\theta G|$, this implies
\eqref{eq:singular-noncollapse}, with the factor two absorbing the scalar
convergence.

For the decay and stable semiseparable masks, Schur's test gives
$\|\bL\|_{\op}\leq C_L$ with $C_L=(1-\gamma)^{-1}$ for decay and
$C_L=(1-a)^{-1}$ for semiseparable masks.  Therefore
\[
    \|\bT_\ell\|_{\op}
    \leq
    C_L\max_i w_i^2.
\]
The maximum of $\ell$ fixed-rank Haar row coordinates satisfies
$\max_iw_i^2=O_{\Prob}(\log\ell)$.  On the other hand, since $L_{ii}=1$,
\[
    \|\bT_\ell\|_F^2
    \geq
    \sum_{i=1}^{\ell}w_i^4,
    \qquad
    \|\bT_\ell\|_*
    \geq
    |\operatorname{Tr}\bT_\ell|
    =
    \sum_{i=1}^{\ell}w_i^2.
\]
Lemma~\ref{lem:empirical-haar-row} gives
$\ell^{-1}\sum_iw_i^4\to\E N(0,1)^4=3$ and
$\ell^{-1}\sum_iw_i^2\to1$.  Thus
\[
    \srank(\bT_\ell)
    =
    \frac{\|\bT_\ell\|_F^2}{\|\bT_\ell\|_{\op}^2}
    \geq
    c\,\frac{\ell}{(\log\ell)^2},
    \qquad
    \frac{\|\bT_\ell\|_*}{\|\bT_\ell\|_{\op}}
    \geq
    c\,\frac{\ell}{\log\ell}
\]
with probability tending to one.  Since Shannon entropy dominates min-entropy,
$\erank_{\mathrm{sing}}(\bT_\ell)\geq
\|\bT_\ell\|_*/\|\bT_\ell\|_{\op}$.  This proves part (ii).  The scalar
prefactor $|\sigma^Q\sigma^K\langle\bv^Q,\bv^K\rangle|/\ell$ does not change
stable rank or entropy effective rank.

For the causal linear mask, the same empirical law gives
\[
    \|\bT_\ell\|_F^2
    =
    \sum_{i\geq j}w_i^2w_j^2
    =
    \frac12\left(\sum_iw_i^2\right)^2
    +
    \frac12\sum_iw_i^4
    =
    \Theta_{\Prob}(\ell^2).
\]
The lower bound $\|\bT_\ell\|_{\op}\geq c\ell$ follows by testing the matrix on
the normalized all-one vector and using
$\ell^{-2}\sum_{i\geq j}w_iw_j\to(\E|N(0,1)|)^2/2=1/\pi$.  Under the weighted
Volterra spread estimate \eqref{eq:weighted-volterra-spread},
$\|\bT_\ell\|_{\op}\leq C_V\ell$, hence
$\srank(\bT_\ell)=\Theta_{\Prob}(1)$.  The same estimate gives
\[
    \erank_{\mathrm{sing}}(\bT_\ell)
    \geq
    \frac{\|\bT_\ell\|_*}{\|\bT_\ell\|_{\op}}
    \geq
    c\log\ell,
\]
which proves part (iii).
\end{proof}

\end{document}
